\documentclass[12pt,a4paper]{article}

\usepackage[utf8]{inputenc}
\usepackage[brazil]{babel}
\usepackage{enumitem}

\title{Questionário Tema 2 – Tomada de Decisão Baseada na Evidência}
\author{Disciplina: Revisão Sistemática e Meta-análises\\Pós-graduação (Mestrado e Doutoramento)}
\date{}

\begin{document}

\maketitle

\section*{Instruções}
Responda às questões abaixo selecionando a(s) opção(ões) correta(s). Cada questão é de múltipla escolha.  

\section*{Questões}

\begin{enumerate}[label=\textbf{Q\arabic*:}]
    
% -------------------------
\item O que caracteriza a \textbf{evidência científica}? 
\begin{enumerate}[label=\alph*)]
    \item É obtida por métodos explícitos e sistemáticos
    \item É sempre absoluta e imutável
    \item É reprodutível ou verificável
    \item Possui incerteza quantificável
\end{enumerate}

% -------------------------
\item Quais são os componentes principais da \textbf{Prática Baseada na Evidência (EBP)}?
\begin{enumerate}[label=\alph*)]
    \item Melhor evidência científica disponível
    \item Experiência profissional especializada
    \item Contexto, valores e necessidades dos utilizadores finais
    \item Intuição pessoal do profissional
\end{enumerate}

% -------------------------
\item Segundo Karl Popper, uma afirmação científica deve ser:
\begin{enumerate}[label=\alph*)]
    \item Falseável
    \item Absoluta
    \item Não revisável
    \item Baseada apenas em opinião
\end{enumerate}

% -------------------------
\item Qual das seguintes opções descreve corretamente o \textbf{fluxo da evidência}?
\begin{enumerate}[label=\alph*)]
    \item Geração → Síntese → Transferência → Implementação → Avaliação
    \item Implementação → Avaliação → Geração → Síntese → Transferência
    \item Síntese → Geração → Avaliação → Transferência → Implementação
    \item Transferência → Implementação → Avaliação → Síntese → Geração
\end{enumerate}

% -------------------------
\item Qual é a função do \textbf{ecossistema da evidência}?
\begin{enumerate}[label=\alph*)]
    \item Integrar múltiplos atores: investigadores, revisores, instituições, decisores e sociedade
    \item Avaliar apenas estudos individuais
    \item Substituir a necessidade de síntese da evidência
    \item Garantir que apenas opiniões de especialistas são consideradas
\end{enumerate}

% -------------------------
\item Por que é necessária a \textbf{síntese da evidência}?
\begin{enumerate}[label=\alph*)]
    \item Para aumentar o poder estatístico
    \item Para reduzir incerteza
    \item Para identificar lacunas de investigação
    \item Para evitar decisões baseadas em estudos isolados
\end{enumerate}

% -------------------------
\item Quais são os tipos de revisão mais comuns?
\begin{enumerate}[label=\alph*)]
    \item Revisão narrativa
    \item Revisão de escopo
    \item Revisão rápida
    \item Revisão sistemática
\end{enumerate}

% -------------------------
\item Quais são os principais elementos de uma revisão sistemática?
\begin{enumerate}[label=\alph*)]
    \item Questão estruturada (PICO)
    \item Protocolo pré-definido
    \item Critérios de inclusão/exclusão
    \item Busca abrangente e avaliação da qualidade
    \item Síntese transparente dos resultados
\end{enumerate}

% -------------------------
\item O que deve ser feito antes de iniciar uma revisão sistemática?
\begin{enumerate}[label=\alph*)]
    \item Verificar revisões existentes
    \item Avaliar viabilidade (tempo, equipa, dados)
    \item Definir objetivos e escopo
    \item Escolher diretrizes metodológicas (PRISMA, Cochrane)
    \item Registar protocolo no PROSPERO
\end{enumerate}

% -------------------------
\item Qual é a sequência correta do \textbf{processo de revisão sistemática}?
\begin{enumerate}[label=\alph*)]
    \item Definição da questão → Elaboração do protocolo → Registo → Busca → Seleção → Extração → Avaliação → Síntese → Reporte
    \item Busca → Seleção → Definição da questão → Extração → Síntese → Avaliação → Reporte → Registo → Protocolo
    \item Definição da questão → Reporte → Busca → Avaliação → Seleção → Extração → Síntese → Registo → Protocolo
    \item Protocolo → Definição da questão → Busca → Seleção → Extração → Avaliação → Síntese → Reporte → Registo
\end{enumerate}

\end{enumerate}

\end{document}
